\subsection*{Raymond FAGES 1911-1983}

Fils du pharmacien de Meyrueis, longtemps Maire et Conseiller Général. Après de
brillantes études, il s’investit dans le scoutisme puis s’engagea au service de la
jeunesse catholique de Lozère à qui il a fait découvrir, après la 2\textsuperscript{o} guerre mondiale
les trésors de la Lozère et ceux du monde.
Aumônier de la JAC (Jeunesse Agricole Catholique), il dirigea bien des sessions de
formation rurale, et créa l’école ménagère de Saint Rome de Dolan. Mais c’est
surtout dans l’Hérault qu’il contribuera à la création de plusieurs maisons au service
de l’enfance handicapée :
L’Institut Médico Educatif (IME) de Batipaume , près d’Agde (appelé désormais
Institut Raymond Fages), un autre IME à Florensac et enfin le château de Bellesague
près de Mende qui lui fut offert par la famille du général de la Porte du Theil, ancien
Directeur des Chantiers de Jeunesse.

\subsection*{Robert HUBAC 1910-1975}

Son père originaire de Conilergues dans la vallée de la Brèze, était pasteur. Agrégé
d’histoire et de géographie, il fut professeur dans les classes de préparation aux
grandes écoles, puis doyen de l’Inspection Générale d’Histoire et de Géographie.
Il pêchait la truite à la mouche, avec adresse aux alentours de sa maison de
Conilergues et me fit remarquer, gentiment, vers 1950, que je pêchais chez lui !
Protestant engagé, il donna bien des articles au journal « Evangile et Liberté », sous
le pseudonyme de Robert Louis (par respect sourilleux de la laïcité).

\subsection*{Mgr MARET 1805-1884}

Né à Meyrueis, où son père était juge de paix. Doyen de la Faculté de Théologie
Catholique de Paris, il prit parti contre l’absolutisme pontifical, pensant qu’il devait
y avoir un équilibre entre le pouvoir du pape et l’autorité des évêques.
Il ne fut pas écouté et lorsque le dogme de l’Infaillibilité du Pape fut promulgué il se
soumit et dû reprendre à son éditeur tous les exemplaires de son ouvrage exposant
ses idées : « Du concile général et de la paix religieuse ». Pour rembourser ses
dettes, il vendit, sur le Causse Méjean, le domaine de la Calvalade, qu’il avait hérité
de son père.

\subsection*{François MARZIALS 1779-1861}

Un des derniers pasteurs cévenols formé au Séminaire de Lausanne créé par Antoine
Court, restaurateur des églises protestantes de France. Il fut pasteur à Barnave dans
la Drôme, puis à Nérac et enfin à Montauban. De tendance « évangélique », il fut élu
président du Consistoire de Montauban, à la suite de son collègue et ami Molines,
libéral, qui avait épousé une « héritière-paysanne » de Carnac, près de Rousses, le
village d’origine de la famille Couderc-Loupiac.
